\title{Highly accurate protein structure prediction with AlphaFold}
\begin{document}
\maketitle

\begin{abstract}
Nature | Vol 596 | 26 August 2021 | 583 Article Highly accurate protein structure prediction with AlphaFold John Jumper1,4 ✉, Richard Evans1,4, Alexander Pritzel1,4, Tim Green1,4, Michael Figurnov1,4, Olaf Ronneberger1,4, Kathryn Tunyasuvunakool1,4, Russ Bates1,4, Augustin Žídek1,4, Anna Potapenko1,4, Alex Bridgland1,4, Clemens Meyer1,4, Simon A. A. Kohl1,4, Andrew J. Ballard1,4, Andrew Cowie1,4, Bernardino Romera-Paredes1,4, Stanislav Nikolov1,4, Rishub Jain1,4, Jonas Adler1, Trevor Back1, Stig Petersen1, David Reiman1, Ellen Clancy1, Michal Zielinski1, Martin Steinegger2,3, Michalina Pacholska1, Tamas Berghammer1, Sebastian Bodenstein1, David Silver1, Oriol Vinyals1, Andrew W. Senior1, Koray Kavukcuoglu1, Pushmeet Kohli1 \& Demis Hassabis1,4 ✉ Proteins are essential to life, and understanding their structure can facilitate a mechanistic understanding of their function. Through an enormous experimental effort1–4, the structures of around 100,000 unique proteins have been determined5, but this represents a small fraction of the billions of known protein sequences6,7. Structural coverage is bottlenecked by the months to years of painstaking effort required to determine a single protein structure. Accurate computational approaches are needed to address this gap and to enable large-scale structural bioinformatics. Predicting the three-dimensional structure that a protein will adopt based solely on its amino acid sequence—the structure prediction component of the ‘protein folding 
\end{abstract}

\section{Recovered Text}
Nature  |  Vol 596  |  26 August 2021  |  583
Article
Highly accurate protein structure prediction
with AlphaFold
John Jumper1,4 ✉, Richard Evans1,4, Alexander Pritzel1,4, Tim Green1,4, Michael Figurnov1,4,
Olaf Ronneberger1,4, Kathryn Tunyasuvunakool1,4, Russ Bates1,4, Augustin Žídek1,4,
Anna Potapenko1,4, Alex Bridgland1,4, Clemens Meyer1,4, Simon A. A. Kohl1,4,
Andrew J. Ballard1,4, Andrew Cowie1,4, Bernardino Romera-Paredes1,4, Stanislav Nikolov1,4,
Rishub Jain1,4, Jonas Adler1, Trevor Back1, Stig Petersen1, David Reiman1, Ellen Clancy1,
Michal Zielinski1, Martin Steinegger2,3, Michalina Pacholska1, Tamas Berghammer1,
Sebastian Bodenstein1, David Silver1, Oriol Vinyals1, Andrew W. Senior1, Koray Kavukcuoglu1,
Pushmeet Kohli1 \& Demis Hassabis1,4 ✉
Proteins are essential to life, and understanding their structure can facilitate a
mechanistic understanding of their function. Through an enormous experimental
effort1–4, the structures of around 100,000 unique proteins have been determined5, but
this represents a small fraction of the billions of known protein sequences6,7. Structural
coverage is bottlenecked by the months to years of painstaking effort required to
determine a single protein structure. Accurate computational approaches are needed
to address this gap and to enable large-scale structural bioinformatics. Predicting the
three-dimensional structure that a protein will adopt based solely on its amino acid
sequence—the structure prediction component of the ‘protein folding problem’8—has
been an important open research problem for more than 50 years9. Despite recent
progress10–14, existing methods fall far short of atomic accuracy, especially when no
homologous structure is available. Here we provide the first computational method
that can regularly predict protein structures with atomic accuracy even in cases in which
no similar structure is known. We validated an entirely redesigned version of our neural
network-based model, AlphaFold, in the challenging 14th Critical Assessment of protein
Structure Prediction (CASP14)15, demonstrating accuracy competitive with
experimental structures in a majority of cases and greatly outperforming other
methods. Underpinning the latest version of AlphaFold is a novel machine learning
approach that incorporates physical and biological knowledge about protein structure,
leveraging multi-sequence alignments, into the design of the deep learning algorithm.
The development of computational methods to predict
three-dimensional (3D) protein structures from the protein sequence
has proceeded along two complementary paths that focus on either the
physical interactions or the evolutionary history. The physical interac-
tion programme heavily integrates our understanding of molecular
driving forces into either thermodynamic or kinetic simulation of pro-
tein physics16 or statistical approximations thereof17. Although theoreti-
cally very appealing, this approach has proved highly challenging for
even moderate-sized proteins due to the computational intractability
of molecular simulation, the context dependence of protein stability
and the difficulty of producing sufficiently accurate models of protein
physics. The evolutionary programme has provided an alternative in
recent years, in which the constraints on protein structure are derived
from bioinformatics analysis of the evolutionary history of proteins,
homology to solved structures18,19 and pairwise evolutionary correla-
tions20–24. This bioinformatics approach has benefited greatly from
the steady growth of experimental protein structures deposited in
the Protein Data Bank (PDB)5, the explosion of genomic sequencing
and the rapid development of deep learning techniques to interpret
these correlations. Despite these advances, contemporary physical
and evolutionary-history-based approaches produce predictions that
are far short of experimental accuracy in the majority of cases in which
a close homologue has not been solved experimentally and this has
limited their utility for many biological applications.
In this study, we develop the first, to our knowledge, computational
approach capable of predicting protein structures to near experimental
accuracy in a majority of cases. The neural network AlphaFold that we
developed was entered into the CASP14 assessment (May–July 2020;
entered under the team name ‘AlphaFold2’ and a completely different
model from our CASP13 AlphaFold system10). The CASP assessment is
carried out biennially using recently solved structures that have not
been deposited in the PDB or publicly disclosed so that it is a blind test
https://doi.org/10.1038/s41586-021-03819-2
Received: 11 May 2021
Accepted: 12 July 2021
Published online: 15 July 2021
Open access
 Check for updates
1DeepMind, London, UK. 2School of Biological Sciences, Seoul National University, Seoul, South Korea. 3Artificial Intelligence Institute, Seoul National University, Seoul, South Korea. 4These
authors contributed equally: John Jumper, Richard Evans, Alexander Pritzel, Tim Green, Michael Figurnov, Olaf Ronneberger, Kathryn Tunyasuvunakool, Russ Bates, Augustin Žídek, Anna
Potapenko, Alex Bridgland, Clemens Meyer, Simon A. A. Kohl, Andrew J. Ballard, Andrew Cowie, Bernardino Romera-Paredes, Stanislav Nikolov, Rishub Jain, Demis Hassabis.
✉e-mail: jumper@deepmind.com; dhcontact@deepmind.com
584  |  Nature  |  Vol 596  |  26 August 2021
Article
for the participating methods, and has long served as the gold-standard
assessment for the accuracy of structure prediction25,26.
In CASP14, AlphaFold structures were vastly more accurate than
competing methods. AlphaFold structures had a median backbone
accuracy of 0.96 Å r.m.s.d.95 (Cα root-mean-square deviation at 95\%
residue coverage) (95\% confidence interval = 0.85–1.16 Å) whereas
the next best performing method had a median backbone accuracy
of 2.8 Å r.m.s.d.95 (95\% confidence interval = 2.7–4.0 Å) (measured on
CASP domains; see Fig. 1a for backbone accuracy and Supplementary
Fig. 14 for all-atom accuracy). As a comparison point for this accuracy,
the width of a carbon atom is approximately 1.4 Å. In addition to very
accurate domain structures (Fig. 1b), AlphaFold is able to produce
highly accurate side chains (Fig. 1c) when the backbone is highly accu-
rate and considerably improves over template-based methods even
when strong templates are available. The all-atom accuracy of Alpha-
Fold was 1.5 Å r.m.s.d.95 (95\% confidence interval = 1.2–1.6 Å) compared
with the 3.5 Å r.m.s.d.95 (95\% confidence interval = 3.1–4.2 Å) of the best
alternative method. Our methods are scalable to very long proteins with
accurate domains and domain-packing (see Fig. 1d for the prediction
of a 2,180-residue protein with no structural homologues). Finally, the
model is able to provide precise, per-residue estimates of its reliability
that should enable the confident use of these predictions.
We demonstrate in Fig. 2a that the high accuracy that AlphaFold dem-
onstrated in CASP14 extends to a large sample of recently released PDB
structures; in this dataset, all structures were deposited in the PDB after
our training data cut-off and are analysed as full chains (see Methods,
Supplementary Fig. 15 and Supplementary Table 6 for more details).
Furthermore, we observe high side-chain accuracy when the back-
bone prediction is accurate (Fig. 2b) and we show that our confidence
measure, the predicted local-distance difference test (pLDDT), reliably
predicts the Cα local-distance difference test (lDDT-Cα) accuracy of the
corresponding prediction (Fig. 2c). We also find that the global super-
position metric template modelling score (TM-score)27 can be accu-
rately estimated (Fig. 2d). Overall, these analyses validate that the high
accuracy and reliability of AlphaFold on CASP14 proteins also transfers
to an uncurated collection of recent PDB submissions, as would be
expected (see Supplementary Methods 1.15 and Supplementary Fig. 11
for confirmation that this high accuracy extends to new folds).
The AlphaFold network
AlphaFold greatly improves the accuracy of structure prediction by
incorporating novel neural network architectures and training proce-
dures based on the evolutionary, physical and geometric constraints
of protein structures. In particular, we demonstrate a new architecture
to jointly embed multiple sequence alignments (MSAs) and pairwise
features, a new output representation and associated loss that enable
accurate end-to-end structure prediction, a new equivariant attention
a
G427
AlphaFold
G009
G473
G129
G403
G032
G420
G480
G498
G488
G368
G324
G362
G253
G216
0
1
2
3
4
Median Cα r.m.s.d.95 (Å)
b
C terminus
N terminus
AlphaFold
Experiment
r.m.s.d.95 = 0.8 Å; TM-score = 0.93
c
AlphaFold
Experiment
r.m.s.d. = 0.59 Å within 8 Å of Zn
d
AlphaFold
Experiment
r.m.s.d.95 = 2.2 Å; TM-score = 0.96
MSA
3D structure
Low
confidence
High
confidence
Templates
Input sequence
MSA
representation
(s,r,c)
Evoformer
(48 blocks)
Structure
module
 (8 blocks)
+
+
← Recycling (three times)
Pairing
Pair
representation
(r,r,c)
Pair
representation
(r,r,c)
Genetic
database
search
Structure
database
search
e
Single repr. (r,c)
Fig. 1 | AlphaFold produces highly accurate structures. a, The performance
of AlphaFold on the CASP14 dataset (n = 87 protein domains) relative to the top-
15 entries (out of 146 entries), group numbers correspond to the numbers
assigned to entrants by CASP. Data are median and the 95\% confidence interval
of the median, estimated from 10,000 bootstrap samples. b, Our prediction of
CASP14 target T1049 (PDB 6Y4F, blue) compared with the true (experimental)
structure (green). Four residues in the C terminus of the crystal structure are
B-factor outliers and are not depicted. c, CASP14 target T1056 (PDB 6YJ1).
An example of a well-predicted zinc-binding site (AlphaFold has accurate side
chains even though it does not explicitly predict the zinc ion). d, CASP target
T1044 (PDB 6VR4)—a 2,180-residue single chain—was predicted with correct
domain packing (the prediction was made after CASP using AlphaFold without
intervention). e, Model architecture. Arrows show the information flow among
the various components described in this paper. Array shapes are shown in
parentheses with s, number of sequences (Nseq in the main text); r, number of
residues (Nres in the main text); c, number of channels.
Nature  |  Vol 596  |  26 August 2021  |  585
architecture, use of intermediate losses to achieve iterative refinement
of predictions, masked MSA loss to jointly train with the structure,
learning from unlabelled protein sequences using self-distillation and
self-estimates of accuracy.
The AlphaFold network directly predicts the 3D coordinates of all
heavy atoms for a given protein using the primary amino acid sequence
and aligned sequences of homologues as inputs (Fig. 1e; see Methods
for details of inputs including databases, MSA construction and use of
templates). A description of the most important ideas and components
is provided below. The full network architecture and training procedure
are provided in the Supplementary Methods.
The network comprises two main stages. First, the trunk of the net-
work processes the inputs through repeated layers of a novel neural
network block that we term Evoformer to produce an Nseq × Nres array
(Nseq, number of sequences; Nres, number of residues) that represents
a processed MSA and an Nres × Nres array that represents residue pairs.
The MSA representation is initialized with the raw MSA (although
see Supplementary Methods 1.2.7 for details of handling very deep
MSAs). The Evoformer blocks contain a number of attention-based
and non-attention-based components. We show evidence in ‘Interpret-
ing the neural network’ that a concrete structural hypothesis arises
early within the Evoformer blocks and is continuously refined. The key
innovations in the Evoformer block are new mechanisms to exchange
information within the MSA and pair representations that enable direct
reasoning about the spatial and evolutionary relationships.
The trunk of the network is followed by the structure module that
introduces an explicit 3D structure in the form of a rotation and transla-
tion for each residue of the protein (global rigid body frames). These
representations are initialized in a trivial state with all rotations set to
the identity and all positions set to the origin, but rapidly develop and
refine a highly accurate protein structure with precise atomic details.
Key innovations in this section of the network include breaking the
chain structure to allow simultaneous local refinement of all parts of
the structure, a novel equivariant transformer to allow the network to
implicitly reason about the unrepresented side-chain atoms and a loss
term that places substantial weight on the orientational correctness
of the residues. Both within the structure module and throughout
the whole network, we reinforce the notion of iterative refinement
by repeatedly applying the final loss to outputs and then feeding the
outputs recursively into the same modules. The iterative refinement
using the whole network (which we term ‘recycling’ and is related to
approaches in computer vision28,29) contributes markedly to accuracy
with minor extra training time (see Supplementary Methods 1.8 for
details).
Evoformer
The key principle of the building block of the network—named Evo-
former (Figs. 1e, 3a)—is to view the prediction of protein structures
as a graph inference problem in 3D space in which the edges of the
graph are defined by residues in proximity. The elements of the pair
representation encode information about the relation between the
residues (Fig. 3b). The columns of the MSA representation encode the
individual residues of the input sequence while the rows represent
the sequences in which those residues appear. Within this framework,
we define a number of update operations that are applied in each block
in which the different update operations are applied in series.
The MSA representation updates the pair representation through an
element-wise outer product that is summed over the MSA sequence
dimension. In contrast to previous work30, this operation is applied
within every block rather than once in the network, which enables the
continuous communication from the evolving MSA representation to
the pair representation.
Within the pair representation, there are two different update pat-
terns. Both are inspired by the necessity of consistency of the pair
0–0.5
0.5–1
1–2
2–4
4–8
>8
Full chain Cα r.m.s.d.95 (Å)
0
0.05
0.10
0.15
0.20
0.25
0.30
Fraction of proteins
a
20
40
60
80
100
lDDT-Cα of a residue
0.5
0.6
0.7
0.8
0.9
1.0
Fraction of correct F1 rotamers
b
0
20
40
60
80
100
Average pLDDT on the resolved region
0
20
40
60
80
100
lDDT-Cα
80
90
100
80
90
100
c
0
0.2
0.4
0.6
0.8
1.0
pTM on the resolved region
0
0.2
0.4
0.6
0.8
1.0
TM-score
0.8
0.9
1.0
0.8
0.9
1.0
d
Fig. 2 | Accuracy of AlphaFold on recent PDB structures. The analysed
structures are newer than any structure in the training set. Further filtering is
applied to reduce redundancy (see Methods). a, Histogram of backbone
r.m.s.d. for full chains (Cα r.m.s.d. at 95\% coverage). Error bars are 95\%
confidence intervals (Poisson). This dataset excludes proteins with a template
(identified by hmmsearch) from the training set with more than 40\% sequence
identity covering more than 1\% of the chain (n = 3,144 protein chains). The
overall median is 1.46 Å (95\% confidence interval = 1.40–1.56 Å). Note that this
measure will be highly sensitive to domain packing and domain accuracy; a
high r.m.s.d. is expected for some chains with uncertain packing or packing
errors. b, Correlation between backbone accuracy and side-chain accuracy.
Filtered to structures with any observed side chains and resolution better than
2.5 Å (n = 5,317 protein chains); side chains were further filtered to
B-factor <30 Å2. A rotamer is classified as correct if the predicted torsion angle
is within 40°. Each point aggregates a range of lDDT-Cα, with a bin size of 2 units
above 70 lDDT-Cα and 5 units otherwise. Points correspond to the mean
accuracy; error bars are 95\% confidence intervals (Student t-test) of the mean
on a per-residue basis. c, Confidence score compared to the true accuracy on
chains. Least-squares linear fit lDDT-Cα = 0.997 × pLDDT − 1.17 (Pearson’s
r = 0.76). n = 10,795 protein chains. The shaded region of the linear fit
represents a 95\% confidence interval estimated from 10,000 bootstrap
samples. In the companion paper39, additional quantification of the reliability
of pLDDT as a confidence measure is provided. d, Correlation between pTM
and full chain TM-score. Least-squares linear fit TM-score = 0.98 × pTM + 0.07
(Pearson’s r = 0.85). n = 10,795 protein chains. The shaded region of the linear fit
represents a 95\% confidence interval estimated from 10,000 bootstrap
samples.
586  |  Nature  |  Vol 596  |  26 August 2021
Article
representation—for a pairwise description of amino acids to be represent-
able as a single 3D structure, many constraints must be satisfied including
the triangle inequality on distances. On the basis of this intuition, we
arrange the update operations on the pair representation in terms of
triangles of edges involving three different nodes (Fig. 3c). In particular,
we add an extra logit bias to axial attention31 to include the ‘missing edge’
of the triangle and we define a non-attention update operation ‘triangle
multiplicative update’ that uses two edges to update the missing third
edge (see Supplementary Methods 1.6.5 for details). The triangle multipli-
cative update was developed originally as a more symmetric and cheaper
replacement for the attention, and networks that use only the attention or
multiplicative update are both able to produce high-accuracy structures.
However, the combination of the two updates is more accurate.
We also use a variant of axial attention within the MSA representation.
During the per-sequence attention in the MSA, we project additional
logits from the pair stack to bias the MSA attention. This closes the loop
by providing information flow from the pair representation back into
the MSA representation, ensuring that the overall Evoformer block is
able to fully mix information between the pair and MSA representations
and prepare for structure generation within the structure module.
End-to-end structure prediction
The structure module (Fig. 3d) operates on a concrete 3D backbone
structure using the pair representation and the original sequence row
(single representation) of the MSA representation from the trunk. The
3D backbone structure is represented as Nres independent rotations
and translations, each with respect to the global frame (residue gas)
(Fig. 3e). These rotations and translations—representing the geometry
of the N-Cα-C atoms—prioritize the orientation of the protein back-
bone so that the location of the side chain of each residue is highly
constrained within that frame. Conversely, the peptide bond geometry
is completely unconstrained and the network is observed to frequently
violate the chain constraint during the application of the structure mod-
ule as breaking this constraint enables the local refinement of all parts
of the chain without solving complex loop closure problems. Satisfac-
tion of the peptide bond geometry is encouraged during fine-tuning
by a violation loss term. Exact enforcement of peptide bond geometry
is only achieved in the post-prediction relaxation of the structure by
gradient descent in the Amber32 force field. Empirically, this final relaxa-
tion does not improve the accuracy of the model as measured by the
f
Backbone frames
(r, 3×3) and (r,3)
Backbone frames
(r, 3×3) and (r,3)
(initially all at the origin)
Predict relative
rotations and
translations
Predict F angles
and compute all
atom positions
(Rk, tk)
xi
8 blocks (shared weights)
d
Pair
representation
(r,r,c)
Single repr. (r,c)
Single repr. (r,c)
MSA
representation
(s,r,c)
Pair
representation
(r,r,c)
Row-wise
gated
self-
attention
with pair
bias
Column-
wise
gated
self-
attention
+
Tran-
sition
Outer
product
mean
+
+
Triangle
update
using
outgoing
edges
Triangle
self-
attention
around
starting
node
Triangle
self-
attention
around
ending
node
Triangle
update
using
incoming
edges
Tran-
sition
MSA
representation
(s,r,c)
Pair
representation
(r,r,c)
a
+
+
+
+
+
+
Triangle self-attention around
starting node
Triangle self-attention around
ending node
ki
ij
kj
jk
ij
ik
i
k
j
i
j
k
ki
ij
kj
i
j
k
ik
ij
jk
i
j
k
Triangle multiplicative update
using ‘outgoing’ edges
Triangle multiplicative update
using ‘incoming’ edges
c
IPA
module
+
Pair representation
(r,r,c)
ij
i
j
k
ik
i
j
k
jk
kj
ji
ki
i
j
k
ij
ji
jk
kj
ki
ik
Corresponding edges
in a graph
b
48 blocks (no shared weights)
e
(Rk, tk)
~
~
xi
~
Fig. 3 | Architectural details. a, Evoformer block. Arrows show the information
flow. The shape of the arrays is shown in parentheses. b, The pair representation
interpreted as directed edges in a graph. c, Triangle multiplicative update and
triangle self-attention. The circles represent residues. Entries in the pair
representation are illustrated as directed edges and in each diagram, the edge
being updated is ij. d, Structure module including Invariant point attention (IPA)
module. The single representation is a copy of the first row of the MSA
representation. e, Residue gas: a representation of each residue as one
free-floating rigid body for the backbone (blue triangles) and χ angles for the
side chains (green circles). The corresponding atomic structure is shown below.
f, Frame aligned point error (FAPE). Green, predicted structure; grey, true
structure; (Rk, tk), frames; xi, atom positions.
Nature  |  Vol 596  |  26 August 2021  |  587
global distance test (GDT)33 or lDDT-Cα34 but does remove distracting
stereochemical violations without the loss of accuracy.
The residue gas representation is updated iteratively in two stages
(Fig. 3d). First, a geometry-aware attention operation that we term
‘invariant point attention’ (IPA) is used to update an Nres set of neural
activations (single representation) without changing the 3D positions,
then an equivariant update operation is performed on the residue gas
using the updated activations. The IPA augments each of the usual
attention queries, keys and values with 3D points that are produced
in the local frame of each residue such that the final value is invariant
to global rotations and translations (see Methods ‘IPA’ for details). The
3D queries and keys also impose a strong spatial/locality bias on the
attention, which is well-suited to the iterative refinement of the protein
structure. After each attention operation and element-wise transition
block, the module computes an update to the rotation and translation
of each backbone frame. The application of these updates within the
local frame of each residue makes the overall attention and update
block an equivariant operation on the residue gas.
Predictions of side-chain χ angles as well as the final, per-residue
accuracy of the structure (pLDDT) are computed with small per-residue
networks on the final activations at the end of the network. The estimate
of the TM-score (pTM) is obtained from a pairwise error prediction that
is computed as a linear projection from the final pair representation. The
final loss (which we term the frame-aligned point error (FAPE) (Fig. 3f))
compares the predicted atom positions to the true positions under
many different alignments. For each alignment, defined by aligning
the predicted frame (Rk, tk) to the corresponding true frame, we com-
pute the distance of all predicted atom positions xi from the true atom
positions. The resulting Nframes × Natoms distances are penalized with a
clamped L1 loss. This creates a strong bias for atoms to be correct relative
to the local frame of each residue and hence correct with respect to its
side-chain interactions, as well as providing the main source of chirality
for AlphaFold (Supplementary Methods 1.9.3 and Supplementary Fig. 9).
Training with labelled and unlabelled data
The AlphaFold architecture is able to train to high accuracy using only
supervised learning on PDB data, but we are able to enhance accuracy
(Fig. 4a) using an approach similar to noisy student self-distillation35.
In this procedure, we use a trained network to predict the structure of
around 350,000 diverse sequences from Uniclust3036 and make a new
dataset of predicted structures filtered to a high-confidence subset. We
then train the same architecture again from scratch using a mixture of
PDB data and this new dataset of predicted structures as the training
data, in which the various training data augmentations such as crop-
ping and MSA subsampling make it challenging for the network to
recapitulate the previously predicted structures. This self-distillation
procedure makes effective use of the unlabelled sequence data and
considerably improves the accuracy of the resulting network.
Additionally, we randomly mask out or mutate individual residues
within the MSA and have a Bidirectional Encoder Representations from
Transformers (BERT)-style37 objective to predict the masked elements of
the MSA sequences. This objective encourages the network to learn to
interpret phylogenetic and covariation relationships without hardcoding
a particular correlation statistic into the features. The BERT objective is
trained jointly with the normal PDB structure loss on the same training
examples and is not pre-trained, in contrast to recent independent work38.
Interpreting the neural network
To understand how AlphaFold predicts protein structure, we trained
a separate structure module for each of the 48 Evoformer blocks in
the network while keeping all parameters of the main network fro-
zen (Supplementary Methods 1.14). Including our recycling stages,
this provides a trajectory of 192 intermediate structures—one per full
Evoformer block—in which each intermediate represents the belief of
the network of the most likely structure at that block. The resulting
trajectories are surprisingly smooth after the first few blocks, show-
ing that AlphaFold makes constant incremental improvements to the
structure until it can no longer improve (see Fig. 4b for a trajectory of
accuracy). These trajectories also illustrate the role of network depth.
For very challenging proteins such as ORF8 of SARS-CoV-2 (T1064),
the network searches and rearranges secondary structure elements
for many layers before settling on a good structure. For other proteins
such as LmrP (T1024), the network finds the final structure within the
first few layers. Structure trajectories of CASP14 targets T1024, T1044,
T1064 and T1091 that demonstrate a clear iterative building process
for a range of protein sizes and difficulties are shown in Supplementary
Videos 1–4. In Supplementary Methods 1.16 and Supplementary Figs. 12,
13, we interpret the attention maps produced by AlphaFold layers.
Figure 4a contains detailed ablations of the components of AlphaFold
that demonstrate that a variety of different mechanisms contribute
to AlphaFold accuracy. Detailed descriptions of each ablation model,
their training details, extended discussion of ablation results and the
–20
–10
0
GDT difference compared
with baseline
No IPA and no recycling
No end-to-end structure gradients
(keep auxiliary heads)
No triangles, biasing or gating
(use axial attention)
No recycling
No auxiliary masked MSA head
No IPA (use direct projection)
No raw MSA
(use MSA pairwise frequencies)
No auxiliary distogram head
No templates
Baseline
With self-distillation training
Test set of CASP14 domains
–4
–2
0
2
lDDT-Cα difference
compared with baseline
Test set of PDB chains
a
0
48
96
144
192
Evoformer block
0
20
40
60
80
100
Domain GDT
T1024 D1
T1024 D2
T1064 D1
b
Fig. 4 | Interpreting the neural network. a, Ablation results on two target sets:
the CASP14 set of domains (n = 87 protein domains) and the PDB test set of
chains with template coverage of ≤30\% at 30\% identity (n = 2,261 protein
chains). Domains are scored with GDT and chains are scored with lDDT-Cα. The
ablations are reported as a difference compared with the average of the three
baseline seeds. Means (points) and 95\% bootstrap percentile intervals (error
bars) are computed using bootstrap estimates of 10,000 samples. b, Domain
GDT trajectory over 4 recycling iterations and 48 Evoformer blocks on CASP14
targets LmrP (T1024) and Orf8 (T1064) where D1 and D2 refer to the individual
domains as defined by the CASP assessment. Both T1024 domains obtain the
correct structure early in the network, whereas the structure of T1064 changes
multiple times and requires nearly the full depth of the network to reach the
final structure. Note, 48 Evoformer blocks comprise one recycling iteration.
588  |  Nature  |  Vol 596  |  26 August 2021
Article
effect of MSA depth on each ablation are provided in Supplementary
Methods 1.13 and Supplementary Fig. 10.
MSA depth and cross-chain contacts
Although AlphaFold has a high accuracy across the vast majority of
deposited PDB structures, we note that there are still factors that affect
accuracy or limit the applicability of the model. The model uses MSAs
and the accuracy decreases substantially when the median alignment
depth is less than around 30 sequences (see Fig. 5a for details). We
observe a threshold effect where improvements in MSA depth over
around 100 sequences lead to small gains. We hypothesize that the MSA
information is needed to coarsely find the correct structure within the
early stages of the network, but refinement of that prediction into a
high-accuracy model does not depend crucially on the MSA information.
The other substantial limitation that we have observed is that AlphaFold
is much weaker for proteins that have few intra-chain or homotypic con-
tacts compared to the number of heterotypic contacts (further details
are provided in a companion paper39). This typically occurs for bridging
domains within larger complexes in which the shape of the protein is
created almost entirely by interactions with other chains in the complex.
Conversely, AlphaFold is often able to give high-accuracy predictions for
homomers, even when the chains are substantially intertwined (Fig. 5b).
We expect that the ideas of AlphaFold are readily applicable to predicting
full hetero-complexes in a future system and that this will remove the dif-
ficulty with protein chains that have a large number of hetero-contacts.
Related work
The prediction of protein structures has had a long and varied develop-
ment, which is extensively covered in a number of reviews14,40–43. Despite
the long history of applying neural networks to structure prediction14,42,43,
they have only recently come to improve structure prediction10,11,44,45.
These approaches effectively leverage the rapid improvement in com-
puter vision systems46 by treating the problem of protein structure
prediction as converting an ‘image’ of evolutionary couplings22–24 to an
‘image’ of the protein distance matrix and then integrating the distance
predictions into a heuristic system that produces the final 3D coordinate
prediction. A few recent studies have been developed to predict the 3D
coordinates directly47–50, but the accuracy of these approaches does not
match traditional, hand-crafted structure prediction pipelines51. In paral-
lel, the success of attention-based networks for language processing52
and, more recently, computer vision31,53 has inspired the exploration of
attention-based methods for interpreting protein sequences54–56.
Discussion
The methodology that we have taken in designing AlphaFold is a combi-
nation of the bioinformatics and physical approaches: we use a physical
and geometric inductive bias to build components that learn from PDB
data with minimal imposition of handcrafted features (for example,
AlphaFold builds hydrogen bonds effectively without a hydrogen bond
score function). This results in a network that learns far more efficiently
from the limited data in the PDB but is able to cope with the complexity
and variety of structural data.
In particular, AlphaFold is able to handle missing the physical context
and produce accurate models in challenging cases such as intertwined
homomers or proteins that only fold in the presence of an unknown
haem group. The ability to handle underspecified structural conditions
is essential to learning from PDB structures as the PDB represents the
full range of conditions in which structures have been solved. In gen-
eral, AlphaFold is trained to produce the protein structure most likely
to appear as part of a PDB structure. For example, in cases in which a
particular stochiometry, ligand or ion is predictable from the sequence
alone, AlphaFold is likely to produce a structure that respects those
constraints implicitly.
AlphaFold has already demonstrated its utility to the experimental
community, both for molecular replacement57 and for interpreting
cryogenic electron microscopy maps58. Moreover, because AlphaFold
outputs protein coordinates directly, AlphaFold produces predictions
in graphics processing unit (GPU) minutes to GPU hours depending on
the length of the protein sequence (for example, around one GPU min-
ute per model for 384 residues; see Methods for details). This opens up
the exciting possibility of predicting structures at the proteome-scale
and beyond—in a companion paper39, we demonstrate the application
of AlphaFold to the entire human proteome39.
The explosion in available genomic sequencing techniques and data
has revolutionized bioinformatics but the intrinsic challenge of experi-
mental structure determination has prevented a similar expansion in
our structural knowledge. By developing an accurate protein structure
100
101
102
103
104
Median per-residue Neff for the chain
0
20
40
60
80
100
lDDT-Cα
Coverage < 0.3
Coverage > 0.6
a
AlphaFold Experiment
b
Fig. 5 | Effect of MSA depth and cross-chain contacts. a, Backbone accuracy
(lDDT-Cα) for the redundancy-reduced set of the PDB after our training data
cut-off, restricting to proteins in which at most 25\% of the long-range contacts
are between different heteromer chains. We further consider two groups of
proteins based on template coverage at 30\% sequence identity: covering more
than 60\% of the chain (n = 6,743 protein chains) and covering less than 30\% of
the chain (n = 1,596 protein chains). MSA depth is computed by counting the
number of non-gap residues for each position in the MSA (using the Neff
weighting scheme; see Methods for details) and taking the median across
residues. The curves are obtained through Gaussian kernel average smoothing
(window size is 0.2 units in log10(Neff)); the shaded area is the 95\% confidence
interval estimated using bootstrap of 10,000 samples. b, An intertwined
homotrimer (PDB 6SK0) is correctly predicted without input stoichiometry
and only a weak template (blue is predicted and green is experimental).
Nature  |  Vol 596  |  26 August 2021  |  589
prediction algorithm, coupled with existing large and well-curated
structure and sequence databases assembled by the experimental
community, we hope to accelerate the advancement of structural
bioinformatics that can keep pace with the genomics revolution. We
hope that AlphaFold—and computational approaches that apply its
techniques for other biophysical problems—will become essential
tools of modern biology.
Online content
Any methods, additional references, Nature Research reporting sum-
maries, source data, extended data, supplementary information,
acknowledgements, peer review information; details of author con-
tributions and competing interests; and statements of data and code
availability are available at https://doi.org/10.1038/s41586-021-03819-2.
1.
Thompson, M. C., Yeates, T. O. \& Rodriguez, J. A. Advances in methods for atomic
resolution macromolecular structure determination. F1000Res. 9, 667 (2020).
2.
Bai, X.-C., McMullan, G. \& Scheres, S. H. W. How cryo-EM is revolutionizing structural
biology. Trends Biochem. Sci. 40, 49–57 (2015).
3.
Jaskolski, M., Dauter, Z. \& Wlodawer, A. A brief history of macromolecular crystallography,
illustrated by a family tree and its Nobel fruits. FEBS J. 281, 3985–4009 (2014).
4.
Wüthrich, K. The way to NMR structures of proteins. Nat. Struct. Biol. 8, 923–925 (2001).
5.
wwPDB Consortium. Protein Data Bank: the single global archive for 3D macromolecular
structure data. Nucleic Acids Res. 47, D520–D528 (2018).
6.
Mitchell, A. L. et al. MGnify: the microbiome analysis resource in 2020. Nucleic Acids Res.
48, D570–D578 (2020).
7.
Steinegger, M., Mirdita, M. \& Söding, J. Protein-level assembly increases protein sequence
recovery from metagenomic samples manyfold. Nat. Methods 16, 603–606 (2019).
8.
Dill, K. A., Ozkan, S. B., Shell, M. S. \& Weikl, T. R. The protein folding problem. Annu. Rev.
Biophys. 37, 289–316 (2008).
9.
Anfinsen, C. B. Principles that govern the folding of protein chains. Science 181, 223–230
(1973).
10.
Senior, A. W. et al. Improved protein structure prediction using potentials from deep
learning. Nature 577, 706–710 (2020).
11.
Wang, S., Sun, S., Li, Z., Zhang, R. \& Xu, J. Accurate de novo prediction of protein contact
map by ultra-deep learning model. PLOS Comput. Biol. 13, e1005324 (2017).
12.
Zheng, W. et al. Deep-learning contact-map guided protein structure prediction in
CASP13. Proteins 87, 1149–1164 (2019).
13.
Abriata, L. A., Tamò, G. E. \& Dal Peraro, M. A further leap of improvement in tertiary
structure prediction in CASP13 prompts new routes for future assessments. Proteins 87,
1100–1112 (2019).
14.
Pearce, R. \& Zhang, Y. Deep learning techniques have significantly impacted protein
structure prediction and protein design. Curr. Opin. Struct. Biol. 68, 194–207 (2021).
15.
Moult, J., Fidelis, K., Kryshtafovych, A., Schwede, T. \& Topf, M. Critical assessment of
techniques for protein structure prediction, fourteenth round. CASP 14 Abstract Book
https://www.predictioncenter.org/casp14/doc/CASP14\_Abstracts.pdf (2020).
16.
Brini, E., Simmerling, C. \& Dill, K. Protein storytelling through physics. Science 370,
eaaz3041 (2020).
17.
Sippl, M. J. Calculation of conformational ensembles from potentials of mean force.
An approach to the knowledge-based prediction of local structures in globular proteins.
J. Mol. Biol. 213, 859–883 (1990).
18.
Šali, A. \& Blundell, T. L. Comparative protein modelling by satisfaction of spatial
restraints. J. Mol. Biol. 234, 779–815 (1993).
19.
Roy, A., Kucukural, A. \& Zhang, Y. I-TASSER: a unified platform for automated protein
structure and function prediction. Nat. Protocols 5, 725–738 (2010).
20.	 Altschuh, D., Lesk, A. M., Bloomer, A. C. \& Klug, A. Correlation of co-ordinated amino acid
substitutions with function in viruses related to tobacco mosaic virus. J. Mol. Biol. 193,
693–707 (1987).
21.
Shindyalov, I. N., Kolchanov, N. A. \& Sander, C. Can three-dimensional contacts in protein
structures be predicted by analysis of correlated mutations? Protein Eng. 7, 349–358 (1994).
22.	 Weigt, M., White, R. A., Szurmant, H., Hoch, J. A. \& Hwa, T. Identification of direct residue
contacts in protein–protein interaction by message passing. Proc. Natl Acad. Sci. USA
106, 67–72 (2009).
23.	 Marks, D. S. et al. Protein 3D structure computed from evolutionary sequence variation.
PLoS ONE 6, e28766 (2011).
24.	 Jones, D. T., Buchan, D. W. A., Cozzetto, D. \& Pontil, M. PSICOV: precise structural contact
prediction using sparse inverse covariance estimation on large multiple sequence
alignments. Bioinformatics 28, 184–190 (2012).
25.	 Moult, J., Pedersen, J. T., Judson, R. \& Fidelis, K. A large-scale experiment to assess protein
structure prediction methods. Proteins 23, ii–iv (1995).
26.	 Kryshtafovych, A., Schwede, T., Topf, M., Fidelis, K. \& Moult, J. Critical assessment of
methods of protein structure prediction (CASP)-round XIII. Proteins 87, 1011–1020 (2019).
27.
Zhang, Y. \& Skolnick, J. Scoring function for automated assessment of protein structure
template quality. Proteins 57, 702–710 (2004).
28.	 Tu, Z. \& Bai, X. Auto-context and its application to high-level vision tasks and 3D brain
image segmentation. IEEE Trans. Pattern Anal. Mach. Intell. 32, 1744–1757 (2010).
29.	 Carreira, J., Agrawal, P., Fragkiadaki, K. \& Malik, J. Human pose estimation with iterative
error feedback. In Proc. IEEE Conference on Computer Vision and Pattern Recognition
4733–4742 (2016).
30.	 Mirabello, C. \& Wallner, B. rawMSA: end-to-end deep learning using raw multiple
sequence alignments. PLoS ONE 14, e0220182 (2019).
31.
Huang, Z. et al. CCNet: criss-cross attention for semantic segmentation. In Proc. IEEE/CVF
International Conference on Computer Vision 603–612 (2019).
32.	 Hornak, V. et al. Comparison of multiple Amber force fields and development of
improved protein backbone parameters. Proteins 65, 712–725 (2006).
33.	 Zemla, A. LGA: a method for finding 3D similarities in protein structures. Nucleic Acids
Res. 31, 3370–3374 (2003).
34.	 Mariani, V., Biasini, M., Barbato, A. \& Schwede, T. lDDT: a local superposition-free score for
comparing protein structures and models using distance difference tests. Bioinformatics
29, 2722–2728 (2013).
35.	 Xie, Q., Luong, M.-T., Hovy, E. \& Le, Q. V. Self-training with noisy student improves
imagenet classification. In Proc. IEEE/CVF Conference on Computer Vision and Pattern
Recognition 10687–10698 (2020).
36.	 Mirdita, M. et al. Uniclust databases of clustered and deeply annotated protein
sequences and alignments. Nucleic Acids Res. 45, D170–D176 (2017).
37.
Devlin, J., Chang, M.-W., Lee, K. \& Toutanova, K. BERT: pre-training of deep bidirectional
transformers for language understanding. In Proc. 2019 Conference of the North
American Chapter of the Association for Computational Linguistics: Human Language
Technologies 1, 4171–4186 (2019).
38.	 Rao, R. et al. MSA transformer. In Proc. 38th International Conference on Machine
Learning PMLR 139, 8844–8856 (2021).
39.	 Tunyasuvunakool, K. et al. Highly accurate protein structure prediction for the human
proteome. Nature https://doi.org/10.1038/s41586-021-03828-1 (2021).
40.	 Kuhlman, B. \& Bradley, P. Advances in protein structure prediction and design. Nat. Rev.
Mol. Cell Biol. 20, 681–697 (2019).
41.
Marks, D. S., Hopf, T. A. \& Sander, C. Protein structure prediction from sequence variation.
Nat. Biotechnol. 30, 1072–1080 (2012).
42.	 Qian, N. \& Sejnowski, T. J. Predicting the secondary structure of globular proteins using
neural network models. J. Mol. Biol. 202, 865–884 (1988).
43.	 Fariselli, P., Olmea, O., Valencia, A. \& Casadio, R. Prediction of contact maps with neural
networks and correlated mutations. Protein Eng. 14, 835–843 (2001).
44.	 Yang, J. et al. Improved protein structure prediction using predicted interresidue
orientations. Proc. Natl Acad. Sci. USA 117, 1496–1503 (2020).
45.	 Li, Y. et al. Deducing high-accuracy protein contact-maps from a triplet of coevolutionary
matrices through deep residual convolutional networks. PLOS Comput. Biol. 17,
e1008865 (2021).
46.	 He, K., Zhang, X., Ren, S. \& Sun, J. Deep residual learning for image recognition. In
Proc. IEEE Conference on Computer Vision and Pattern Recognition 770–778 (2016).
47.
AlQuraishi, M. End-to-end differentiable learning of protein structure. Cell Syst. 8,
292–301 (2019).
48.	 Senior, A. W. et al. Protein structure prediction using multiple deep neural networks in the
13th Critical Assessment of Protein Structure Prediction (CASP13). Proteins 87, 1141–1148
(2019).
49.	 Ingraham, J., Riesselman, A. J., Sander, C. \& Marks, D. S. Learning protein structure with a
differentiable simulator. in Proc. International Conference on Learning Representations
(2019).
50.	 Li, J. Universal transforming geometric network. Preprint at https://arxiv.org/
abs/1908.00723 (2019).
51.
Xu, J., McPartlon, M. \& Li, J. Improved protein structure prediction by deep learning
irrespective of co-evolution information. Nat. Mach. Intell. 3, 601–609 (2021).
52.	 Vaswani, A. et al. Attention is all you need. In Advances in Neural Information Processing
Systems 5998–6008 (2017).
53.	 Wang, H. et al. Axial-deeplab: stand-alone axial-attention for panoptic segmentation. in
European Conference on Computer Vision 108–126 (Springer, 2020).
54.	 Alley, E. C., Khimulya, G., Biswas, S., AlQuraishi, M. \& Church, G. M. Unified rational
protein engineering with sequence-based deep representation learning. Nat. Methods 16,
1315–1322 (2019).
55.	 Heinzinger, M. et al. Modeling aspects of the language of life through transfer-learning
protein sequences. BMC Bioinformatics 20, 723 (2019).
56.	 Rives, A. et al. Biological structure and function emerge from scaling unsupervised
learning to 250 million protein sequences. Proc. Natl Acad. Sci. USA 118, e2016239118
(2021).
57.
Pereira, J. et al. High-accuracy protein structure prediction in CASP14. Proteins https://doi.
org/10.1002/prot.26171 (2021).
58.	 Gupta, M. et al. CryoEM and AI reveal a structure of SARS-CoV-2 Nsp2, a multifunctional
protein involved in key host processes. Preprint at https://doi.org/10.1101/2021.05.10.
443524 (2021).
Publisher’s note Springer Nature remains neutral with regard to jurisdictional claims in
published maps and institutional affiliations.
Open Access This article is licensed under a Creative Commons Attribution
4.0 International License, which permits use, sharing, adaptation, distribution
and reproduction in any medium or format, as long as you give appropriate
credit to the original author(s) and the source, provide a link to the Creative Commons license,
and indicate if changes were made. The images or other third party material in this article are
included in the article’s Creative Commons license, unless indicated otherwise in a credit line
to the material. If material is not included in the article’s Creative Commons license and your
intended use is not permitted by statutory regulation or exceeds the permitted use, you will
need to obtain permission directly from the copyright holder. To view a copy of this license,
visit http://creativecommons.org/licenses/by/4.0/.
© The Author(s) 2021
Article
Methods
Full algorithm details
Extensive explanations of the components and their motivations are
available in Supplementary Methods 1.1–1.10, in addition, pseudocode
is available in Supplementary Information Algorithms 1–32, network
diagrams in Supplementary Figs. 1–8, input features in Supplementary
Table 1 and additional details are provided in Supplementary Tables 2, 3.
Training and inference details are provided in Supplementary Methods
1.11–1.12 and Supplementary Tables 4, 5.
IPA
The IPA module combines the pair representation, the single repre-
sentation and the geometric representation to update the single rep-
resentation (Supplementary Fig. 8). Each of these representations
contributes affinities to the shared attention weights and then uses
these weights to map its values to the output. The IPA operates in 3D
space. Each residue produces query points, key points and value points
in its local frame. These points are projected into the global frame using
the backbone frame of the residue in which they interact with each
other. The resulting points are then projected back into the local frame.
The affinity computation in the 3D space uses squared distances and
the coordinate transformations ensure the invariance of this module
with respect to the global frame (see Supplementary Methods 1.8.2
‘Invariant point attention (IPA)’ for the algorithm, proof of invariance
and a description of the full multi-head version). A related construc-
tion that uses classic geometric invariants to construct pairwise fea-
tures in place of the learned 3D points has been applied to protein
design59.
In addition to the IPA, standard dot product attention is computed
on the abstract single representation and a special attention on the pair
representation. The pair representation augments both the logits and
the values of the attention process, which is the primary way in which
the pair representation controls the structure generation.
Inputs and data sources
Inputs to the network are the primary sequence, sequences from evo-
lutionarily related proteins in the form of a MSA created by standard
tools including jackhmmer60 and HHBlits61, and 3D atom coordinates
of a small number of homologous structures (templates) where avail-
able. For both the MSA and templates, the search processes are tuned
for high recall; spurious matches will probably appear in the raw MSA
but this matches the training condition of the network.
One of the sequence databases used, Big Fantastic Database (BFD),
was custom-made and released publicly (see ‘Data availability’) and
was used by several CASP teams. BFD is one of the largest publicly avail-
able collections of protein families. It consists of 65,983,866 families
represented as MSAs and hidden Markov models (HMMs) covering
2,204,359,010 protein sequences from reference databases, metage-
nomes and metatranscriptomes.
BFD was built in three steps. First, 2,423,213,294 protein sequences
were collected from UniProt (Swiss-Prot\&TrEMBL, 2017-11)62, a soil refer-
ence protein catalogue and the marine eukaryotic reference catalogue7,
and clustered to 30\% sequence identity, while enforcing a 90\% align-
ment coverage of the shorter sequences using MMseqs2/Linclust63.
This resulted in 345,159,030 clusters. For computational efficiency,
we removed all clusters with less than three members, resulting in
61,083,719 clusters. Second, we added 166,510,624 representative pro-
tein sequences from Metaclust NR (2017-05; discarding all sequences
shorter than 150 residues)63 by aligning them against the cluster rep-
resentatives using MMseqs264. Sequences that fulfilled the sequence
identity and coverage criteria were assigned to the best scoring cluster.
The remaining 25,347,429 sequences that could not be assigned were
clustered separately and added as new clusters, resulting in the final
clustering. Third, for each of the clusters, we computed an MSA using
FAMSA65 and computed the HMMs following the Uniclust HH-suite
database protocol36.
The following versions of public datasets were used in this study. Our
models were trained on a copy of the PDB5 downloaded on 28 August
2019. For finding template structures at prediction time, we used a copy
of the PDB downloaded on 14 May 2020, and the PDB7066 clustering
database downloaded on 13 May 2020. For MSA lookup at both training
and prediction time, we used Uniref9067 v.2020\_01, BFD, Uniclust3036
v.2018\_08 and MGnify6 v.2018\_12. For sequence distillation, we used
Uniclust3036 v.2018\_08 to construct a distillation structure dataset.
Full details are provided in Supplementary Methods 1.2.
For MSA search on BFD + Uniclust30, and template search against
PDB70, we used HHBlits61 and HHSearch66 from hh-suite v.3.0-beta.3
(version 14/07/2017). For MSA search on Uniref90 and clustered MGnify,
we used jackhmmer from HMMER368. For constrained relaxation of
structures, we used OpenMM v.7.3.169 with the Amber99sb force field32.
For neural network construction, running and other analyses, we used
TensorFlow70, Sonnet71, NumPy72, Python73 and Colab74.
To quantify the effect of the different sequence data sources, we
re-ran the CASP14 proteins using the same models but varying how the
MSA was constructed. Removing BFD reduced the mean accuracy by
0.4 GDT, removing Mgnify reduced the mean accuracy by 0.7 GDT, and
removing both reduced the mean accuracy by 6.1 GDT. In each case, we
found that most targets had very small changes in accuracy but a few
outliers had very large (20+ GDT) differences. This is consistent with the
results in Fig. 5a in which the depth of the MSA is relatively unimportant
until it approaches a threshold value of around 30 sequences when the
MSA size effects become quite large. We observe mostly overlapping
effects between inclusion of BFD and Mgnify, but having at least one
of these metagenomics databases is very important for target classes
that are poorly represented in UniRef, and having both was necessary
to achieve full CASP accuracy.
Training regimen
To train, we use structures from the PDB with a maximum release date
of 30 April 2018. Chains are sampled in inverse proportion to cluster
size of a 40\% sequence identity clustering. We then randomly crop
them to 256 residues and assemble into batches of size 128. We train the
model on Tensor Processing Unit (TPU) v3 with a batch size of 1 per TPU
core, hence the model uses 128 TPU v3 cores. The model is trained until
convergence (around 10 million samples) and further fine-tuned using
longer crops of 384 residues, larger MSA stack and reduced learning
rate (see Supplementary Methods 1.11 for the exact configuration). The
initial training stage takes approximately 1 week, and the fine-tuning
stage takes approximately 4 additional days.
The network is supervised by the FAPE loss and a number of auxil-
iary losses. First, the final pair representation is linearly projected to
a binned distance distribution (distogram) prediction, scored with
a cross-entropy loss. Second, we use random masking on the input
MSAs and require the network to reconstruct the masked regions
from the output MSA representation using a BERT-like loss37. Third,
the output single representations of the structure module are used to
predict binned per-residue lDDT-Cα values. Finally, we use an auxiliary
side-chain loss during training, and an auxiliary structure violation loss
during fine-tuning. Detailed descriptions and weighting are provided
in the Supplementary Information.
An initial model trained with the above objectives was used to make
structure predictions for a Uniclust dataset of 355,993 sequences with
the full MSAs. These predictions were then used to train a final model
with identical hyperparameters, except for sampling examples 75\% of
the time from the Uniclust prediction set, with sub-sampled MSAs, and
25\% of the time from the clustered PDB set.
We train five different models using different random seeds, some
with templates and some without, to encourage diversity in the predic-
tions (see Supplementary Table 5 and Supplementary Methods 1.12.1
for details). We also fine-tuned these models after CASP14 to add a
pTM prediction objective (Supplementary Methods 1.9.7) and use the
obtained models for Fig. 2d.
Inference regimen
We inference the five trained models and use the predicted confidence
score to select the best model per target.
Using our CASP14 configuration for AlphaFold, the trunk of the net-
work is run multiple times with different random choices for the MSA
cluster centres (see Supplementary Methods 1.11.2 for details of the
ensembling procedure). The full time to make a structure prediction
varies considerably depending on the length of the protein. Repre-
sentative timings for the neural network using a single model on V100
GPU are 4.8 min with 256 residues, 9.2 min with 384 residues and 18 h
at 2,500 residues. These timings are measured using our open-source
code, and the open-source code is notably faster than the version we
ran in CASP14 as we now use the XLA compiler75.
Since CASP14, we have found that the accuracy of the network with-
out ensembling is very close or equal to the accuracy with ensembling
and we turn off ensembling for most inference. Without ensembling,
the network is 8× faster and the representative timings for a single
model are 0.6 min with 256 residues, 1.1 min with 384 residues and
2.1 h with 2,500 residues.
Inferencing large proteins can easily exceed the memory of a single
GPU. For a V100 with 16 GB of memory, we can predict the structure
of proteins up to around 1,300 residues without ensembling and the
256- and 384-residue inference times are using the memory of a single
GPU. The memory usage is approximately quadratic in the number of
residues, so a 2,500-residue protein involves using unified memory so
that we can greatly exceed the memory of a single V100. In our cloud
setup, a single V100 is used for computation on a 2,500-residue protein
but we requested four GPUs to have sufficient memory.
Searching genetic sequence databases to prepare inputs and final
relaxation of the structures take additional central processing unit
(CPU) time but do not require a GPU or TPU.
Metrics
The predicted structure is compared to the true structure from the
PDB in terms of lDDT metric34, as this metric reports the domain accu-
racy without requiring a domain segmentation of chain structures.
The distances are either computed between all heavy atoms (lDDT)
or only the Cα atoms to measure the backbone accuracy (lDDT-Cα).
As lDDT-Cα only focuses on the Cα atoms, it does not include the pen-
alty for structural violations and clashes. Domain accuracies in CASP
are reported as GDT33 and the TM-score27 is used as a full chain global
superposition metric.
We also report accuracies using the r.m.s.d.95 (Cα r.m.s.d. at 95\% cov-
erage). We perform five iterations of (1) a least-squares alignment of the
predicted structure and the PDB structure on the currently chosen Cα
atoms (using all Cα atoms in the first iteration); (2) selecting the 95\%
of Cα atoms with the lowest alignment error. The r.m.s.d. of the atoms
chosen for the final iterations is the r.m.s.d.95. This metric is more robust
to apparent errors that can originate from crystal structure artefacts,
although in some cases the removed 5\% of residues will contain genuine
modelling errors.
Test set of recent PDB sequences
For evaluation on recent PDB sequences (Figs. 2a–d, 4a, 5a), we used
a copy of the PDB downloaded 15 February 2021. Structures were fil-
tered to those with a release date after 30 April 2018 (the date limit for
inclusion in the training set for AlphaFold). Chains were further filtered
to remove sequences that consisted of a single amino acid as well as
sequences with an ambiguous chemical component at any residue
position. Exact duplicates were removed, with the chain with the most
resolved Cα atoms used as the representative sequence. Subsequently,
structures with less than 16 resolved residues, with unknown residues
or solved by NMR methods were removed. As the PDB contains many
near-duplicate sequences, the chain with the highest resolution was
selected from each cluster in the PDB 40\% sequence clustering of the
data. Furthermore, we removed all sequences for which fewer than
80 amino acids had the alpha carbon resolved and removed chains with
more than 1,400 residues. The final dataset contained 10,795 protein
sequences.
The procedure for filtering the recent PDB dataset based on prior
template identity was as follows. Hmmsearch was run with default
parameters against a copy of the PDB SEQRES fasta downloaded
15 February 2021. Template hits were accepted if the associated struc-
ture had a release date earlier than 30 April 2018. Each residue position
in a query sequence was assigned the maximum identity of any template
hit covering that position. Filtering then proceeded as described in
the individual figure legends, based on a combination of maximum
identity and sequence coverage.
The MSA depth analysis was based on computing the normalized
number of effective sequences (Neff) for each position of a query
sequence. Per-residue Neff values were obtained by counting the num-
ber of non-gap residues in the MSA for this position and weighting the
sequences using the Neff scheme76 with a threshold of 80\% sequence
identity measured on the region that is non-gap in either sequence.
Reporting summary
Further information on research design is available in the Nature
Research Reporting Summary linked to this paper.
Data availability
All input data are freely available from public sources.
Structures from the PDB were used for training and a
\end{document}
